% from ega4-2-fr.tex
\oldpage[IV-2]{190}

$\operatorname{codim}(\overline{\{y'\}}\cap\overline{\{x'\}},
\overline{\{x'\}})\geqslant2$.
Comme $y'$ est arbitraire dans
$f^{-1}(T)\cap\overline{\{x'\}}$, on a bien
$\operatorname{codim}(f^{-1}(T)\cap\overline{\{x'\}},
\overline{\{x'\}})\geqslant2$ dans $X'$. C.Q.F.D.
\end{proof}

\begin{theorem}[7.2.5]
\label{IV.2.s7.2.5-fr}
Soit $A$ un anneau local noethérien. Les conditions suivantes sont
équivalentes:
\begin{enumerate}
\item[a)] Pour tout anneau quotient intègre $B$ de $A$, le complété
$\widehat B$ est strictement équidimensionnel (7.2.1).
\item[b)] $A$ est formellement caténaire (7.1.9) et les fibres du morphisme
canonique $\operatorname{Spec}(\widehat A)\to\operatorname{Spec}(A)$
vérifient $(\mathrm S_1)$ (autrement dit, n'ont pas de cycle premier associé
immergé).
\item[c)] $A$ est universellement caténaire (5.6.2) et pour tout anneau
quotient intègre $B$ de $A$, l'anneau $B^{(1)}$ est une $B$-algèbre finie
(cf. (7.2.4)).
\item[d)] $A$ est universellement caténaire et pour tout anneau quotient
intègre $B$ de $A$ et tout anneau local $C=B_{\mathfrak q}$ en un idéal
premier $\mathfrak q$ de $B$, tel que $\dim(C)\geqslant2$, l'anneau
$C^{(\omega)}$ est une $C$-algèbre finie.
\item[e)] $A$ est universellement caténaire et pour tout anneau quotient
intègre $B$ de $A$ et tout anneau local $C=B_{\mathfrak q}$ en un idéal
premier $\mathfrak q$ de $B$, tel que $\dim(C)\geqslant2$, l'anneau complété
$\widehat C$ est tel que $\operatorname{Spec}(\widehat C)$ n'ait pas de
cycle premier associé de dimension 1.
\end{enumerate}
De plus, lorsque ces conditions sont remplies, alors, pour tout anneau
quotient $B$ de $A$ qui est strictement équidimensionnel, le complété
$\widehat B$ est strictement équidimensionnel.
\end{theorem}

\begin{proof}
L'équivalence de c), d) et e) a été démontrée dans (7.2.4). Pour montrer
que a) et b) sont équivalentes, rappelons (7.1.9) que dire que $A$ est
formellement caténaire signifie que pour tout anneau quotient intègre $B$
de $A$, $\widehat B$ est équidimensionnel. D'autre part, toute fibre du
morphisme $\operatorname{Spec}(\widehat A)\to\operatorname{Spec}(A)$ en un
point $x\in\operatorname{Spec}(A)$ est la fibre du morphisme
$\operatorname{Spec}(\widehat B)\to\operatorname{Spec}(B)$ au point
générique $x$ de $\operatorname{Spec}(B)$, où
$B=A/\mathfrak J_x$; dire que cette fibre est sans cycle premier associé
immergé équivaut à dire que $\widehat B$ n'a pas d'idéaux premiers associés
immergés, en vertu de (3.3.3); ceci prouve l'équivalence de a) et b). Le
même raisonnement montre que si a) est vérifiée, alors, pour tout anneau
quotient $B$ de $A$ qui est sans idéaux premiers associés immergés, le
complété $\widehat B$ est aussi sans idéaux premiers associés immergés;
d'autre part, l'hypothèse que $A$ est formellement caténaire entraîne que
si $B$ est équidimensionnel, il en est de même de $\widehat B$ (7.1.8);
ceci établit la dernière assertion du théorème.

Montrons maintenant que a) implique c). La condition a) implique que $A$
est universellement caténaire (7.1.11); montrons d'autre part que pour tout
anneau quotient intègre $B$ de $A$, $B^{(1)}$ est alors une $B$-algèbre
finie. Compte tenu de (7.2.3) appliqué à $B$, il s'agit de montrer que si
$X=\operatorname{Spec}(B)$, si $T$ est une partie fermée de $X$, de
codimension $\geqslant2$, $X'=\operatorname{Spec}(\widehat B)$, et
$g:X'\to X$ le morphisme canonique, alors on a, pour tout
$x'\in\operatorname{Ass}(\mathscr O_{X'})$,
$\operatorname{codim}(g^{-1}(T)\cap\overline{\{x'\}},
\overline{\{x'\}})\geqslant2$.
Mais par hypothèse $X'$ n'a pas de cycle premier associé immergé, donc
$\inf(\operatorname{codim}(g^{-1}(T)\cap\overline{\{x'\}},
\overline{\{x'\}}))$, lorsque $x'$ parcourt
$\operatorname{Ass}(\mathscr O_{X'})$, n'est autre que
$\operatorname{codim}(g^{-1}(T),X')$.
Or, comme $g$ est un morphisme fidèlement plat, on a (6.1.4)
\[
\operatorname{codim}(g^{-1}(T),X')=\operatorname{codim}(T,X)\geqslant2.
\]

